% Chapter 5: Pronouns: The Noun Replacers

\chapter{Pronouns: The Noun Replacers}

\section{Pedagogical Purpose}
After learning to identify and classify nouns, students need to understand how to refer to them efficiently. Pronouns are the primary tool for this. This chapter is crucial for developing fluency and avoiding the awkward repetition of nouns in both speech and writing.

\begin{learningobjectives}
The learner can:
\begin{itemize}
    \item define what a pronoun is and why it is used.
    \item identify personal pronouns (I, we, you, he, she, it, they, me, us, him, her, them).
    \item distinguish between subject pronouns and object pronouns.
    \item replace nouns in sentences with appropriate pronouns.
    \item use pronouns correctly to write clear and concise sentences.
\end{itemize}
\end{learningobjectives}

\section{Prerequisite Check}
\begin{quickcheck}
Underline the nouns in the sentences below.
\begin{enumerate}
    \item \underline{Rohan} kicked the \underline{ball}.
    \item \underline{Maya} gave a \underline{book} to her \underline{teacher}.
    \item The \underline{monkeys} ate the \underline{bananas}.
\end{enumerate}
\end{quickcheck}

\section{Language Experience (Let's Begin)}
Ms. Sharma is talking to Maya and Rohan in the classroom.

\textbf{Ms. Sharma:} "Rohan, can you tell me about your morning?"

\textbf{Rohan:} "Rohan woke up. Then Rohan brushed Rohan's teeth. Rohan's mother gave Rohan breakfast. Rohan liked the breakfast."

\textbf{Maya:} (Giggles) "That sounds funny! You said 'Rohan' so many times."

\textbf{Ms. Sharma:} "Exactly, Maya! It sounds a bit strange, doesn't it? Rohan, how would you normally say that?"

\textbf{Rohan:} "I would say: \textbf{I} woke up. Then \textbf{I} brushed \textbf{my} teeth. \textbf{My} mother gave \textbf{me} breakfast. \textbf{I} liked \textbf{it}."

\textbf{Ms. Sharma:} "Perfect! See how you replaced the noun 'Rohan' with words like 'I' and 'me'? And you replaced 'breakfast' with 'it'. These replacing words are called \textbf{Pronouns}. They make our sentences sound much more natural and less repetitive."

\section{Concept Discovery (What is a Pronoun?)}
A \textbf{pronoun} is a word that is used in place of a noun.

Using pronouns helps us avoid repeating the same nouns over and over again.

\begin{itemize}
    \item \textbf{Without Pronoun:} \textbf{Maya} has a new book. \textbf{Maya} is reading the book.
    \item \textbf{With Pronoun:} \textbf{Maya} has a new book. \textbf{She} is reading \textbf{it}.
\end{itemize}

Here, \texttt{She} replaces \texttt{Maya}, and \texttt{it} replaces \texttt{the book}.

\section{Concept Explanation (Personal Pronouns)}
The most common pronouns are \textbf{personal pronouns}. They change their form depending on their job in a sentence. They can be a \textbf{subject} (the one doing the action) or an \textbf{object} (the one receiving the action).

\subsection*{Subject Pronouns}
\textbf{Subject pronouns} perform the action in a sentence. They usually come before the verb.

\begin{examplebox}
    \textbf{He} threw the ball. (\texttt{He} is the one doing the throwing.)
\end{examplebox}

\subsection*{Object Pronouns}
\textbf{Object pronouns} receive the action in a sentence. They usually come after the verb.

\begin{examplebox}
    Rohan threw the ball to \textbf{him}. (\texttt{him} is the one receiving the ball.)
\end{examplebox}

\subsection*{The Pronoun Team}
Let's meet the whole team of personal pronouns!

\begin{table}[h!]
\centering
\begin{tabular}{|l|l|l|}
\hline
\textbf{Person} & \textbf{Singular} & \textbf{Plural} \\
\hline
\textbf{First Person} & \textbf{I} (Subject) & \textbf{we} (Subject) \\
(the person speaking) & \textbf{me} (Object) & \textbf{us} (Object) \\
\hline
\textbf{Second Person} & \textbf{you} (Subject) & \textbf{you} (Subject) \\
(the person spoken to) & \textbf{you} (Object) & \textbf{you} (Object) \\
\hline
\textbf{Third Person} & \textbf{he, she, it} (Subject) & \textbf{they} (Subject) \\
(the person/thing spoken about) & \textbf{him, her, it} (Object) & \textbf{them} (Object) \\
\hline
\end{tabular}
\caption{Personal Pronouns}
\end{table}

\section{Worked Example}
\textbf{Ms. Sharma:} "Let's practice replacing nouns with the correct pronouns."

\textbf{Original Sentence:} \texttt{Rohan loves cricket. Rohan plays cricket every day.}

\textbf{Let's Think Step-by-Step:}
\begin{enumerate}
    \item \textbf{Identify the noun being repeated:} \texttt{Rohan}.
    \item \textbf{What pronoun can replace \texttt{Rohan}?} \texttt{Rohan} is a single boy being spoken about (Third Person, Singular, Masculine). The pronoun is \texttt{he}.
    \item \textbf{Is \texttt{Rohan} the subject or object in the second sentence?} He is the one playing, so he is the subject. We need the subject pronoun.
    \item \textbf{Replace the noun:} \texttt{He}.
\end{enumerate}
\textbf{Revised Sentence:} \texttt{Rohan loves cricket. \textbf{He} plays cricket every day.}

\section{Guided Practice (Let's Practise Together)}
\begin{activitybox}{Activity A: Choose the Correct Pronoun}
\textbf{Ms. Sharma:} "Let's choose the right pronoun for each sentence. Remember to check if it's the subject or the object. I'll do the first one."
\begin{enumerate}
    \item \textbf{(She / Her)} is my best friend. \textit{Answer: She}
    \item The teacher gave the book to \textbf{(he / him)}.
    \item \textbf{(They / Them)} are playing in the garden.
    \item Rohan is waiting for \textbf{(we / us)}.
\end{enumerate}
\end{activitybox}

\section{Independent Practice (Now, Your Turn!)}
\begin{activitybox}{Activity B: Fill in the Blanks}
Complete the sentences with the correct pronoun from the box.

\fbox{I, me, he, him, she, her, it, we, us, they, them}

\begin{enumerate}
    \item My name is Rohan. ______ am ten years old.
    \item That is Maya. Please give this book to ______.
    \item The dog is hungry. Can you give ______ some food?
    \item My friends are here. ______ are waiting outside.
    \item Rohan and I are going to the park. Do you want to come with ______?
\end{enumerate}
\end{activitybox}

\section{Summary}
\begin{summarybox}
    \begin{itemize}
        \item \textbf{Pronouns} are words that take the place of nouns.
        \item \textbf{Subject Pronouns} do the action (\texttt{I, you, he, she, it, we, they}).
        \item \textbf{Object Pronouns} receive the action (\texttt{me, you, him, her, it, us, them}).
    \end{itemize}
\end{summarybox}